\section{Introduction}

Segmentation metrics summarize different properties of a prediction. Mean
intersection over union (mIoU) measures region overlap; Boundary IoU emphasizes
contours; tail statistics emphasize poorly covered components; and
false-positive scores penalize unsupported predictions. A score that separates
failed from non-failed images need not rank models in the order a declared
policy or operational utility would. Treating these uses as interchangeable
obscures the decision a metric is meant to support, and the problem is sharper
outside the general-scene benchmarks on which most scores are developed.

We study the issue through pairwise preferences over complete six-model image
groups (Fig.~\ref{fig:teaser}). For each image and model pair, seven criteria
may agree, tie, or disagree. The minority count among non-tied preferences
lower-bounds the disagreement of any single ordering. Across eight panels in
three domains the bound is \NBAllMin{}--\NBAllMax{}, and an exhaustive search
over all 720 strict orderings of each image shows it is essentially attained:
the exact minimum exceeds it by at most \NBGapAllMax{}. In social-choice terms,
this exact minimum is the Kemeny--Young optimal disagreement of the criteria
viewed as voters~\cite{kemeny1959mathematics}, a quantitative complement to
Arrow's impossibility theorem~\cite{arrow1951social}. No scalar score, however
it is built, can order these models without contradicting this share of the
declared preferences.

Given an unavoidable conflict, we evaluate a simple protocol: select one score
formula per endpoint on COCO and VOC, freeze the formulas, and score each
endpoint with its own formula once the endpoint is declared. Against reuse of a
single mixed-failure score, this improves rank agreement and severe-miss
detection on all eight panels, including recent ADE20K architectures, street
scenes, and CT, none of which was used for selection.

A single failure score behaves differently. Selected on general scenes, the
mixed score is the best failure detector there, but on street scenes and CT it
trails a single coverage or tail metric. Those domains have a different failure
composition: high-error failures are rare and nearly every failure is a severe
miss, often a small distant pedestrian in street scenes. Re-selecting the
failure score on Cityscapes training data, whose failures share that
composition, restores it on three held-out panels, including CT.

Our contributions are:
(1)~an empirical, essentially exact bound on scalar-ranking disagreement among
declared segmentation criteria, computed as the Kemeny--Young optimum of the
criteria viewed as voters, on eight panels in three domains;
(2)~a cross-domain evaluation of frozen endpoint-specific scoring, which
improves every non-trivial endpoint on every panel; and
(3)~evidence that a single failure score transfers only between domains with
similar failure composition, and that selecting it on data with matching
composition restores it.
All results are conditional on the declared criteria, panels, and endpoints;
they support naming the decision endpoint and its target failure composition
when reporting segmentation quality, not a new universal metric.
